\begin{abstract}
Emerging advancements in research, exemplified by BitNet~\cite{bitnet}, signal the advent of a new generation of 1-bit Large Language Models (LLMs). Here, we present a novel 1-bit LLM architecture termed \textbf{\text{BitNet b1.58}}, in which each parameter (or weight) is strictly ternary, selected from \{-1, 0, 1\}. This model achieves comparable performance in both perplexity and downstream tasks relative to standard full-precision Transformer LLMs (using FP16 or BF16), provided equivalent model scale and training data. Additionally, it offers substantial efficiency improvements with respect to latency, memory usage, computational throughput, and energy consumption. Notably, the 1.58-bit LLM introduces a fundamentally new scaling law and provides a systematic approach for training next-generation LLMs that are both high-performing and efficient. It also lays the foundation for a novel computational framework, encouraging the development of specialized hardware tailored for 1-bit LLM architectures.
\end{abstract}

\section{The Era of 1-bit LLMs}

Recent advances in artificial intelligence have accelerated the expansion of both the scale and effectiveness of Large Language Models (LLMs). Although these models have attained exceptional results across diverse natural language processing tasks, their escalating parameter counts introduce substantial difficulties in terms of deployment. Moreover, concerns regarding their resource consumption and environmental footprint have intensified due to high energy demands.

One promising method for mitigating these limitations involves the application of post-training quantization, where models are compressed to use fewer bits for inference~\cite{smoothquant, gptq, quip, quip_sharp}. By decreasing the precision of both weights and activations, quantization can drastically lower memory footprint and computational expense. While the focus has progressively shifted from 16-bit models toward even lower-precision formats, such as 4-bit quantization~\cite{gptq, awq}, post-training quantization remains a sub-optimal solution—despite its widespread adoption in industrial settings.

Recent advancements in 1-bit model architectures—exemplified by BitNet~\cite{bitnet}—suggest a compelling pathway for lowering LLM operating costs without compromising model quality. Standard LLMs employ 16-bit floating-point representations (such as FP16 or BF16), with matrix multiplication serving as the computational backbone. As a consequence, most of the computational overhead arises from floating-point multiplication and addition. BitNet, on the other hand, executes matrix multiplications solely through integer addition, resulting in substantial reductions in energy usage for LLMs. Since power constraints often govern achievable compute performance on modern hardware, these power savings may also facilitate accelerated computation.

Beyond raw computation, model parameter transfer from DRAM to on-chip accelerator memory (such as SRAM) introduces notable expenses during inference. While scaling up SRAM can mitigate bandwidth limitations, it dramatically raises cost relative to DRAM. In comparison to conventional full-precision models, 1-bit LLMs impose dramatically lower demands on both memory capacity and memory bandwidth. Consequently, transferring weights from DRAM becomes less time-consuming and less costly, thereby improving inference speed and overall efficiency.

This study proposes an enhanced 1-bit LLM, termed \textbf{\text{BitNet b1.58}}, in which all parameters are restricted to ternary values from the set \{-1, 0, 1\}. Introducing the additional value 0 to the conventional 1-bit BitNet framework extends the bit representation to 1.58 bits per parameter. \text{BitNet b1.58} maintains all foundational advantages of the original BitNet, including a matrix multiplication approach that is virtually free of multiplications and readily optimized. The model matches the original BitNet in terms of power efficiency, while exceeding full-precision (FP16) LLMs in memory usage, throughput, and latency. Moreover, \text{BitNet b1.58} confers two further benefits: First, its expressive capacity is improved via explicit feature selection facilitated by the inclusion of 0, which meaningfully enhances the potential performance of 1-bit LLMs. Second, empirical evaluation reveals that \text{BitNet b1.58} achieves parity with FP16 baselines on both perplexity and downstream tasks—beginning from the 3B parameter scale—when all other model configurations (such as size and training token counts) are held constant.

\section{BitNet b1.58}

\text{BitNet b1.58} builds upon the BitNet framework, a variant of Transformer that substitutes \emph{nn.Linear} layers with \emph{BitLinear} layers. This model is trained from the beginning, employing weights quantized to 1.58 bits and activations to 8 bits. Relative to the original BitNet, several modifications have been made, as outlined below.

\paragraph{Quantization Function.} In order to limit the weight values to the discrete set \{-1, 0, +1\}, we use an \emph{absmean} quantization strategy. This method initially normalizes the weight matrix by dividing each entry by the mean of its absolute values, after which it rounds each value to its nearest neighbor in \{-1, 0, +1\}:

\begin{equation}
    \widetilde{W} = \mathrm{RoundClip}(\frac{W}{\gamma + \epsilon}, -1, 1),
\end{equation}

\begin{equation}
    \mathrm{RoundClip}(x, a, b) = \max(a, \min(b, \mathrm{round}(x))),
\end{equation}

\begin{equation}
    \gamma = \frac{1}{nm}\sum_{ij} |W_{ij}|.
\end{equation}

For activations, the quantization procedure mirrors that of BitNet, but with a key difference: we do not rescale activations prior to non-linear transformations into the $[0, Q_b]$ interval. Instead, for each token, activations are rescaled to lie within $[-Q_b, Q_b]$, thereby eliminating the need for zero-point quantization. This approach streamlines both the implementation and the system-level optimization, and in our experiments, does not noticeably affect model performance.

\paragraph{LLaMA-alike Components.}

LLaMA~\cite{llama, llama2} has emerged as the standard reference architecture for open-source LLMs. To ensure compatibility with the broader open-source ecosystem, the design of \text{BitNet b1.58} integrates several LLaMA-alike components. In particular, it incorporates RMSNorm~\cite{rmsnorm}, SwiGLU~\cite{swiglu}, and rotary embedding~\cite{rope}, and omits all bias terms. As a result, \text{BitNet b1.58} can be readily deployed within widely used open-source platforms (such as Huggingface, vLLM~\cite{vllm}, and llama.cpp\footnote{\url{https://github.com/ggerganov/llama.cpp}}) with little additional effort.

\section{Results}

We conducted a comparative analysis of \text{BitNet b1.58} and our FP16-reproduced \text{LLaMA LLM} models across multiple scales. To maintain experimental parity, all models underwent pre-training on the RedPajama dataset~\cite{redpajama} for a total of 100 billion tokens. The assessment of zero-shot capabilities was performed over a diverse suite of language benchmarks, such as ARC-Easy~\cite{arc}, ARC-Challenge~\cite{arc}, Hellaswag~\cite{hellaswag}, Winogrande~\cite{winoGrande}, PIQA~\cite{piqa}, OpenbookQA~\cite{openbookqa}, and BoolQ~\cite{boolq}. Additionally, we computed the validation perplexity scores using the WikiText2~\cite{wikitext2} and C4~\cite{c4} corpora.

For both \text{LLaMA LLM} and \text{BitNet b1.58}, we measured GPU memory usage and inference latency during runtime. These metrics were gathered with the FasterTransformer\footnote{\url{https://github.com/NVIDIA/FasterTransformer}} library, recognized for its high-performance inference capabilities on GPU platforms. For \text{BitNet b1.58}, we incorporated the 2-bit kernel from Ladder~\cite{ladder}. The principal inference cost, namely, the time per output token, was the main metric reported.

In Table~\ref{tab:ppl}, we present a summary of perplexity as well as the computational cost associated with both \text{BitNet b1.58} and \text{LLaMA LLM}. The findings indicate that, with model sizes reaching 3B parameters, \text{BitNet b1.58} attains comparable perplexity to the FP16 \text{LLaMA LLM}, while achieving a 2.71$\times$ increase in speed and requiring 3.55$\times$ less GPU memory. Notably, \text{BitNet b1.58} at 3.9B parameters demonstrates a 2.4$\times$ acceleration, a 3.32$\times$ memory reduction, and yet surpasses the \text{LLaMA LLM} 3B in performance.

Table~\ref{tab:zero-shot} provides detailed zero-shot accuracy results for downstream language tasks. Evaluation was conducted using the \emph{lm-evaluation-harness}\footnote{\url{https://github.com/EleutherAI/lm-evaluation-harness}} framework. The results underscore a diminishing performance disparity between \text{BitNet b1.58} and \text{LLaMA LLM} as model size increases, with \text{BitNet b1.58} matching the full-precision baseline from the 3B size onward. As also observed with perplexity outcomes, the downstream accuracy results show that \text{BitNet b1.58} at 3.9B not only achieves superior performance to \text{LLaMA LLM} 3B, but does so at reduced memory and latency costs. This highlights \text{BitNet b1.58} as delivering a Pareto-optimal advancement relative to state-of-the-art LLM models.

\paragraph{Memory and Latency}

We extended our experiments by increasing the model size to 7B, 13B, and 70B parameters and assessed the corresponding computational cost. As depicted in Figure~\ref{fig:latency-memory-energy}, both latency and memory utilization display distinct scaling patterns; notably, performance gains in terms of speed-up are amplified as the model size increases. For instance, \text{BitNet b1.58} at 70B achieves a 4.1-fold speed-up relative to the \text{LLaMA LLM} baseline. This enhanced efficiency arises because the computational overhead associated with \emph{nn.Linear} layers grows in proportion to the model's size. Memory requirements also follow this scaling, but since the embedding layer is always stored at full precision, its share of total memory consumption diminishes in larger models. The latency and memory figures presented were obtained using a 2-bit kernel, implying further optimizations may lower resource costs even more.

\paragraph{Energy}

To analyze energy efficiency, we estimated the arithmetic energy demands for both \text{BitNet b1.58} and \text{LLaMA LLM}. Emphasis was placed on matrix multiplication operations, as these dominate the total energy profile in LLM inference. Figure~\ref{fig:energy} provides a breakdown of energy consumption. For \text{BitNet b1.58}, the majority of operations are INT8 additions, whereas \text{LLaMA LLM} heavily relies on FP16 additions and multiplications. Citing the energy modeling from~\cite{energycost, pokebnn}, matrix multiplications performed by \text{BitNet b1.58} use 71.4 times less energy compared to their counterparts on 7nm hardware. Furthermore, we report end-to-end energy usage for processing 512 tokens per model. The data indicates that, with increasing model size, \text{BitNet b1.58} delivers superior energy efficiency relative to the FP16 \text{LLaMA LLM} reference. This advantage stems from the rising contribution of \emph{nn.Linear} layers in larger models, while the energy proportion associated with other elements decreases.

\paragraph{Throughput}

We assess the throughput of \text{BitNet b1.58} alongside the \text{LLaMA LLM} with 70B parameters, utilizing two A100 GPUs equipped with 80GB memory. To accommodate the \text{LLaMA LLM} 70B model, pipeline parallelism~\cite{gpipe} is employed. The batch size is progressively increased until the GPU memory capacity is fully utilized, maintaining a sequence length of 512. As detailed in Table~\ref{tab:throughput}, \text{BitNet b1.58} 70B supports batch sizes up to 11 times greater than those feasible with \text{LLaMA LLM}, translating to an 8.9-fold improvement in throughput.

\textbf{\text{BitNet b1.58} demonstrates a novel scaling relationship between model performance and inference expenditure.} For comparative purposes, equivalences among model sizes for 1.58-bit and 16-bit precision are summarized, referencing Figure~\ref{fig:latency-memory-energy} and \ref{fig:energy}.

\begin{itemize}
    \item The 13B BitNet b1.58 outperforms the 3B FP16 LLM regarding latency, memory efficiency, and energy utilization.
    \item The 30B BitNet b1.58 surpasses the 7B FP16 LLM in terms of latency, memory demands, and energy consumption.
    \item The 70B BitNet b1.58 exceeds the 13B FP16 LLM on latency, memory footprint, and energy usage.
\end{itemize}

\paragraph{Training with 2T Tokens}

Training token count significantly affects LLM capabilities. To investigate the token-level scaling properties of \text{BitNet b1.58}, we trained this model variant on 2T tokens using the StableLM-3B~\cite{StableLM-3B-4E1T} data regimen, which represents the leading open-source 3B model. Both models were benchmarked against a suite comprising Winogrande~\cite{winoGrande}, PIQA~\cite{piqa}, SciQ~\cite{sciq}, LAMBADA~\cite{lambada}, and ARC-easy~\cite{arc}. Zero-shot accuracy metrics are presented in Table~\ref{tab:2t}. For datasets reporting both standard and normalized accuracy, their mean is used. Performance values for StableLM 3b with 2T tokens are extracted directly from its official report. The evaluation reveals that \text{BitNet b1.58} consistently outperforms on all tested tasks, underscoring the robust generalization potential of 1.58-bit LLMs.

\section{Discussion and Future Work}

\textbf{1-bit Mixture-of-Experts (MoE) LLMs}

Mixture-of-Experts (MoE) architectures have established themselves as an efficient paradigm for large language models (LLMs), notably reducing the computational demands measured by FLOPs. Nevertheless, substantial memory requirements and the costs associated with inter-chip communication often impede their practical implementation. Employing 1.58-bit LLMs offers a compelling solution to these issues. The drastically lowered memory requirements can reduce the hardware footprint necessary for MoE deployment, and can also mitigate the communication overhead incurred during activation exchanges across devices. Ideally, if the whole model fits onto a single chip, inter-device communication costs could be eliminated entirely.

\textbf{Native Support for Long Sequences in LLMs}

The capacity to process extended input sequences natively is increasingly essential in LLM applications. A significant bottleneck for inference with long sequences lies in the memory usage due to key-value (KV) caches. \text{BitNet b1.58} advances the state of the art by lowering activation precision from 16 bits to 8 bits, thereby enabling twice the context length within the same memory budget. Future avenues include further compressing activation precision, potentially losslessly, down to 4 bits or even further for 1.58-bit LLMs, which remains a subject for future research.

\textbf{LLMs for Edge and Mobile Devices}

Deploying 1.58-bit LLMs could significantly boost the viability of running language models on edge and mobile platforms, where computational and memory resources are often at a premium. The energy efficiency and reduced storage footprint of 1.58-bit LLMs not only facilitate their deployment in resource-constrained environments but also open doors to applications previously considered infeasible. This advancement can substantially broaden the functional capabilities of edge and mobile devices, ushering in innovative uses of LLMs. Additionally, the compatibility of 1.58-bit LLMs with CPUs—typically the main processors in these settings—allows models like \text{BitNet b1.58} to execute efficiently, thereby enhancing both performance and capability on such hardware.

\textbf{Hardware Innovation for 1-bit LLMs}

The development of specialized hardware, such as in~Groq\footnote{\url{https://groq.com/}}, demonstrates considerable promise in optimizing infrastructure for LLM workloads (e.g., using LPUs). Building upon this direction, we propose and encourage future research and engineering efforts aimed at designing hardware and systems specifically tailored for 1-bit LLMs, leveraging the computational efficiencies pioneered in BitNet~\cite{bitnet}.